\documentclass[a4paper,11pt]{amsart}

\usepackage[margin=1in]{geometry}
\usepackage{amsmath,amssymb}
\usepackage{booktabs}
\usepackage{hyperref}

\hypersetup{colorlinks=true,linkcolor=blue,citecolor=blue,urlcolor=blue}
\setlength{\emergencystretch}{2em}

\title{The shortest rosary on seven symbols}
\date{}

\newtheorem{theorem}{Theorem}
\newtheorem{lemma}[theorem]{Lemma}
\newtheorem{corollary}[theorem]{Corollary}
\theoremstyle{remark}
\newtheorem{remark}[theorem]{Remark}

\newcommand{\word}{\texttt{123456712347654321765432}}

\begin{document}

\begin{abstract}
A rosary of degree $n$ is a cyclic word over $\{1,\dots,n\}$ in which every
permutation of $\{1,\dots,n\}$ can be read as a subsequence by starting at a
suitable bead and going once around the circle.  Gupta conjectured in 1981
that the shortest rosary of degree $n$ has $\lfloor n^2/2\rfloor$ beads; the
conjecture fails for $n=6$ and $n=8$, and the value for $n=7$ was open.  We
prove that the shortest rosary of degree $7$ has $24$ beads.  The upper bound
comes from a known construction; the lower bound is proved by an exhaustive
computer search, based on a reduction that turns the insertion of a new
symbol into a hitting-set problem on the gaps of a smaller rosary.
\end{abstract}

\maketitle

\section{The problem and the answer}

Write $[n]=\{1,\dots,n\}$.  For a word $w=w_0w_1\cdots w_{L-1}$ and
$0\le s<L$, let $w^{(s)}=w_s w_{s+1}\cdots w_{L-1}w_0\cdots w_{s-1}$ be its
rotation starting at position $s$.  Following Gupta \cite{Gupta1981}, a word
$w$ over $[n]$ is a \emph{rosary of degree $n$} if every permutation of $[n]$,
written as a word, is a subsequence of some rotation of $w$.  Equivalently,
the letters of $w$ are beads on a circle, and each permutation can be read by
choosing a starting bead and moving once around the circle in a fixed
direction, skipping beads freely and never passing the starting bead again.
Let $r(n)$ be the least length of a rosary of degree $n$.

Gupta conjectured that $r(n)=\lfloor n^2/2\rfloor$ and reported that this
was verified for $2\le n\le 6$ \cite[p.~273]{Gupta1981}; his example for
$n=6$ has $18$ beads.  The upper-bound form $r(n)\le n^2/2$ is recorded as
problem E22 in \cite{Guy2004}.  Lecouturier and Zmiaikou proved the upper
bound for even $n$, exhibited rosaries of lengths $17$ for $n=6$ and $31$ for
$n=8$ (so the exact conjecture fails for these $n$), and proved
$r(n)<n^2/2+n/4-1$ for odd $n$, which gives $r(7)\le 25$
\cite{LZ2012}; see also the survey \cite{EV2021}.  The values
$r(2),\dots,r(6)=2,4,8,12,17$ are listed in \cite{A373728}.  A rosary of
length $24$ for $n=7$, obtained by shortening the Lecouturier--Zmiaikou
construction, was posted by Ever-Hadani \cite{EverHadani2026}, together with
an unfinished attempt to exclude length $23$.

\begin{theorem}\label{thm:main}
$r(7)=24$.
\end{theorem}

\begin{corollary}
Gupta's conjectured value $\lfloor n^2/2\rfloor$ is correct for $n=7$.
\end{corollary}

The upper bound is the word
\[
  W_{24}=\word ,
\]
that is, the Lecouturier--Zmiaikou word $(1\cdots7)(1\cdots5)(7\cdots1)(7\cdots2)$
for $n=7$ \cite[Theorem~2]{LZ2012} with one bead $5$ removed.  A direct check
of all $5040$ permutations confirms that $W_{24}$ is a rosary.  The rest of
this note proves the lower bound.  The proof is computer assisted: the
reduction in Section~\ref{sec:reduction} is proved by hand, and it reduces
the lower bound to a finite search, described in Section~\ref{sec:search}.
The same search also gives $r(6)\ge 17$, so $r(6)=17$; we found no earlier
published proof of this lower bound.

\section{The reduction}\label{sec:reduction}

All words below are cyclic.  For a word $v$ of length $m$ we index its
\emph{gaps} by $0,\dots,m-1$, gap $g$ being the position immediately before
the bead $v_g$.

\begin{lemma}[Padding]\label{lem:pad}
Inserting a bead anywhere in a rosary of degree $n$ gives a rosary of degree
$n$.  Consequently, if $r(n)\le L$ then there is a rosary of degree $n$ of
length exactly $L$.
\end{lemma}

\begin{proof}
Every rotation of the old word, starting at a bead $b$, is a subsequence of
the rotation of the new word starting at the same bead.
\end{proof}

\begin{lemma}[Deletion]\label{lem:del}
If $w$ is a rosary of degree $n$ and $a\in[n]$, then deleting every bead $a$
from $w$ gives a rosary on the alphabet $[n]\setminus\{a\}$.
\end{lemma}

\begin{proof}
Let $\pi$ be a permutation of $[n]\setminus\{a\}$.  The permutation
$\pi a$ is a subsequence of some rotation of $w$; deleting the beads $a$
from that rotation leaves a rotation of the shorter word that still contains
$\pi$.
\end{proof}

\begin{lemma}[Cyclic shifts]\label{lem:shift}
If $\pi_1\cdots\pi_n$ can be read in one lap of a cyclic word $w$, then so can
every cyclic shift $\pi_j\cdots\pi_n\pi_1\cdots\pi_{j-1}$.
\end{lemma}

\begin{proof}
Unroll the circle and suppose $\pi_1,\dots,\pi_n$ are read at positions
$i_1<\dots<i_n<i_1+L$.  Then
$i_j<\dots<i_n<i_1+L<\dots<i_{j-1}+L<i_j+L$, which reads the shift in one
lap starting at position $i_j$.
\end{proof}

The key step describes when a new symbol can be added.  For a permutation
$P$ of $[k]$ and a cyclic word $v$ over $[k]$ of length $m$, let
\[
  G_v(P)=\{\,g\in\{0,\dots,m-1\}:\ P \text{ is a subsequence of } v^{(g)}\,\}
\]
be the set of gaps from which $P$ can be read in one lap.

\begin{lemma}[Gap criterion]\label{lem:gap}
Let $v$ be a cyclic word over $[k]$, and let $w$ be obtained by inserting
$c\ge1$ beads of a new symbol $x$ into gaps of $v$, occupying the set of gaps
$S$.  Then $w$ is a rosary on $[k]\cup\{x\}$ if and only if $S\cap G_v(P)\ne
\emptyset$ for every permutation $P$ of $[k]$.
\end{lemma}

\begin{proof}
Every permutation of $[k]\cup\{x\}$ is a cyclic shift of exactly one word of
the form $xP$ with $P$ a permutation of $[k]$.  By Lemma~\ref{lem:shift},
$w$ is a rosary if and only if every word $xP$ can be read in one lap of $w$.
If $xP$ is read with its first letter at a bead $x$ lying in gap $g$, then
$P$ is read from the beads that follow within the lap; ignoring the other
beads $x$, these are the beads of $v^{(g)}$ in order, so $g\in G_v(P)$.
Conversely, if $g\in S\cap G_v(P)$, start at a bead $x$ in gap $g$ and read
$P$ from $v^{(g)}$.
\end{proof}

\begin{lemma}[Symmetries]\label{lem:sym}
Rotating, reversing, or relabelling a rosary gives a rosary.
\end{lemma}

\begin{proof}
Rotation and relabelling are clear.  Reversal maps the rotations of $w$ to
the reversed rotations of the reversed word and maps the set of all
permutations to itself.
\end{proof}

\section{The search}\label{sec:search}

By Lemma~\ref{lem:pad}, it suffices to show that no rosary of degree $7$ has
length $23$.  Suppose $w$ is one.  Every symbol occurs, so the frequencies,
sorted as $c_1\le\dots\le c_7$, form one of the $164$ partitions of $23$ into
seven positive parts.  Relabel so that symbol $i$ occurs $c_i$ times.  For
$k=1,\dots,7$ let $w_k$ be $w$ with the symbols $k+1,\dots,7$ deleted.  By
Lemma~\ref{lem:del} each $w_k$ is a rosary of degree $k$, and by
Lemma~\ref{lem:gap} the beads of symbol $k+1$ occupy a set of gaps of $w_k$
meeting $G_{w_k}(P)$ for every permutation $P$ of $[k]$.

For a fixed profile the search therefore builds the sets $U_1,\dots,U_7$,
where $U_1$ is the single word $1^{c_1}$ and $U_{k+1}$ consists of all words
obtained from some $v\in U_k$ by inserting $c_{k+1}$ beads $k+1$ whose
occupied gaps meet every $G_v(P)$.  Words are stored in a canonical form,
the lexicographically least spelling under rotation, reversal, and
relabelling within classes of equal frequency; by Lemma~\ref{lem:sym} this
loses nothing.  A child $u$ of $v$ is kept only if deleting the largest
symbol from the canonical form of $u$ and canonicalising again returns $v$.
This rule is complete by induction on $k$: if $u$ is any rosary at level
$k+1$ with canonical form $c$, the canonical parent $p$ of $c$ is a rosary by
Lemma~\ref{lem:del}, so $p\in U_k$ by induction.  Applying to $c$ the
rotation, reversal and relabelling that carry the deleted word of $c$ to $p$
gives a word $c'$ equivalent to $u$ that arises from $p$ by an insertion, and
$c'$ passes the test.  Inside a level, placements are enumerated by
branching on the smallest occupied gap of an unhit constraint; a branch is
pruned when the number of pairwise disjoint unhit constraints exceeds the
number of beads left.  At the last level only the existence of a placement
is needed.

The search terminates with $U_7=\emptyset$ for all $164$ profiles.  The
numbers of distinct words, up to symmetry, reached at each level over all
profiles are shown below.

\begin{center}
\begin{tabular}{lccccccc}
\toprule
level $k$ & 1 & 2 & 3 & 4 & 5 & 6 & 7\\
\midrule
words in $U_k$ & 3 & 10 & 88 & 3122 & 125334 & 145472 & 0\\
\bottomrule
\end{tabular}
\end{center}

Hence no rosary of degree $7$ has $23$ beads, and with $W_{24}$ this proves
Theorem~\ref{thm:main}.  For $n=6$ the same program finds no rosary of length
$16$ among the $35$ profiles and finds rosaries of length $17$, so $r(6)=17$.

\section{Verification}

The search is load-bearing, so it was carried out several times.  The
original program and a second implementation of the same method both report
$U_7=\emptyset$, with the level counts above.  Three further programs were
written independently of these.  One inserts a most frequent symbol at each
level and keeps a single set of canonical words under the full symmetry
group; it reports the same level counts and, as a positive control, finds
rosaries of length $24$.  A second deletes a most frequent symbol and checks
every candidate directly against all permutations; it finds no rosary of
length $23$ and exactly $24$ classes of rosaries of length $17$ for $n=6$.
A third tests each accepted child directly, with no pruning, and also finds
none.  For $n\le5$ these programs agree with a brute-force check of all
words of the relevant lengths, including all $5^{13}$ words of length $13$,
and a SAT solver confirms that no rosary of degree $6$ has length $16$.
Deliberately introduced errors in the pruning and in the one-lap test were
detected by these comparisons.  The programs, their outputs, and
instructions for rerunning them are available at
\url{https://github.com/DannyExperiments/seven-symbol-rosary}.

All programs rely on Lemmas~\ref{lem:pad}--\ref{lem:sym}, except the brute
force and SAT checks for $n\le6$.  No part of the argument has been
formally verified.

\begin{remark}
Lemma~\ref{lem:del} is the circular form of a deletion argument used for
linear words by Newey \cite{Newey1973}, and Lemma~\ref{lem:gap} is a circular
analogue of a one-letter characterisation of Radomirovi\'c
\cite[Theorem~5]{Radomirovic2012}.  The canonical-parent rule is a form of
canonical augmentation \cite{McKay1998}.
\end{remark}

\section*{AI disclosure}
Under human direction, an OpenAI Codex cloud agent found the proof and wrote
the first two search programs, and Claude (Anthropic) independently audited
and reimplemented the search, carried out the literature review, and prepared
this manuscript.  Neither system is listed as an author.  The work has been
internally reviewed, not peer reviewed.

\renewcommand{\bibliofont}{\footnotesize\raggedright}
\bibliographystyle{amsplain}
\bibliography{references}
\end{document}
